\lysh{15042_SOLUTION}{1}
\hfill \textbf{\textLambda\textUpsilon\textSigma\textEta} \hfill
\vspace{0.5\baselineskip}

\begin{enumerate}[\upshape \bfseries \alpha)]
    \item Αρκεί να αποδειχθεί ότι υπάρχει $\mu \in \mathbb{R}$ έτσι ώστε $\vec{B\Gamma} = \mu \vec{BM}$. \\
    Πράγματι, θεωρώντας το $B$ ως σημείο αναφοράς, είναι:
    \[
    \vec{AB} - 2\vec{AM} + \vec{A\Gamma} = \vec{0} \Rightarrow
    \]
    \[
    \vec{AB} - 2(\vec{BM} - \vec{BA}) + (\vec{B\Gamma} - \vec{BA}) = \vec{0} \Rightarrow
    \]
    \[
    \vec{AB} + 2\vec{BA} + \vec{B\Gamma} - \vec{BA} = 2\vec{BM} \Rightarrow
    \]
    \[
    \vec{AB} + \vec{BA} + \vec{B\Gamma} = 2\vec{BM} \Rightarrow \vec{B\Gamma} = 2\vec{BM}
    \]

    \item Το $M$ είναι το μέσο του τμήματος $B\Gamma$, διότι
    \[
    \vec{B\Gamma} = 2\vec{BM} \Rightarrow \vec{B\Gamma} - \vec{BM} = \vec{BM} \Rightarrow \vec{M\Gamma} = \vec{BM}
    \]

    \item \begin{enumerate}[\upshape \bfseries i)]
        \item Επειδή τα σημεία $A, B, \Gamma$ ως κορυφές τριγώνου δεν είναι συνευθειακά, έπεται ότι τα μη μηδενικά διανύσματα $\vec{AB}, \vec{A\Gamma}$ δεν είναι παράλληλα. \\
        Είναι $\kappa \vec{A\Gamma} = \lambda \vec{AB}$. \\
        Αν $\kappa \neq 0$, τότε $\kappa \vec{A\Gamma} = \lambda \vec{AB} \Rightarrow \vec{A\Gamma} = \frac{\lambda}{\kappa} \vec{AB} \Rightarrow \vec{A\Gamma} // \vec{AB}$. \\
        Αν $\lambda \neq 0$, τότε $\kappa \vec{A\Gamma} = \lambda \vec{AB} \Rightarrow \vec{AB} = \frac{\kappa}{\lambda} \vec{A\Gamma} \Rightarrow \vec{A\Gamma} // \vec{AB}$. \\
        Επομένως πρέπει $\kappa = \lambda = 0$.

        \item Τότε είναι:
        \[
        \begin{cases}
            \vec{AB} \cdot \vec{A\Gamma} = 0 \\
            \vec{AM} \cdot \vec{B\Gamma} = 0
        \end{cases}
        \]
        Αφού $\vec{AB} \cdot \vec{A\Gamma} = 0$ έπεται ότι τα διανύσματα $\vec{AB}, \vec{A\Gamma}$ είναι κάθετα. Επομένως, το τρίγωνο είναι ορθογώνιο, με $\hat{A} = 90^\circ$. \\
        Αφού $\vec{AM} \cdot \vec{B\Gamma} = 0$ έπεται ότι η διάμεσος $AM$ του ορθογώνιου τριγώνου είναι κάθετη στην πλευρά $B\Gamma$, δηλαδή είναι και ύψος. \\
        Ως εκ τούτου το τρίγωνο είναι και ισοσκελές με $AB = A\Gamma$.
    \end{enumerate}
\end{enumerate}
\end{lysh}